\chapter{Duchenne Muscular Dystrophy}

\section*{Definition and Overview}
Duchenne muscular dystrophy (DMD) is a genetic condition that causes progressive weakness and wasting of the muscles. It is caused by a mutation in the gene that produces dystrophin, a protein essential for keeping muscle fibres intact. Without dystrophin, muscle cells are damaged each time they contract, leading to progressive weakness that begins in early childhood. DMD predominantly affects boys, who are usually diagnosed between the ages of three and five. This chapter explains the condition, its progression, and the treatments and support available to help affected boys and men live as well as possible.

\section*{Epidemiology}
Duchenne muscular dystrophy is the most common and severe form of muscular dystrophy, affecting approximately 1 in 3,500 to 5,000 boys born worldwide. Because the faulty gene is on the X chromosome, the condition affects almost exclusively males, while females may be carriers who usually have no symptoms but can pass the condition to their sons. There is no cure, but advances in respiratory and cardiac care have dramatically improved life expectancy, with many affected men now living into their late twenties and thirties.

\section*{Aetiology and Pathophysiology}
Duchenne muscular dystrophy is caused by mutations in the DMD gene, located on the X chromosome, which provides the instructions for making dystrophin.
\begin{itemize}
    \item \textbf{Genetic cause:} the DMD gene mutation leads to the complete absence of functional dystrophin
    \item \textbf{Inheritance:} X-linked recessive inheritance; affected boys inherit the mutation from their carrier mothers, although a significant minority of cases arise from new mutations
    \item \textbf{Mechanism of damage:} dystrophin normally anchors the muscle cell's internal framework to its outer membrane, protecting it during contraction; without it, the membrane tears, muscle cells are destroyed, and are progressively replaced by fat and scar tissue
    \item \textbf{Progression:} muscle loss begins in the hips, thighs, and shoulders, then spreads to other muscles, including the heart and breathing muscles
    \item \textbf{Cardiac involvement:} the heart muscle is affected, leading to cardiomyopathy in the teenage years
\end{itemize}

\section*{Clinical Features}
The signs of DMD usually become apparent in early childhood.
\begin{itemize}
    \item Delayed milestones, such as late walking
    \item A waddling gait and difficulty running, jumping, and climbing stairs
    \item Frequent falls and trouble getting up from the floor, often climbing up their own legs (Gowers' sign)
    \item Enlarged calf muscles, which feel firm and rubbery because of fat and scar tissue
    \item Progressive weakness, with loss of the ability to walk typically between the ages of 8 and 14
    \item Learning and behavioural difficulties in some boys, including delayed speech and attention problems
    \item Breathing weakness in the teenage years, requiring respiratory support
    \item Heart problems, which may be silent until advanced
\end{itemize}

\section*{Differential Diagnosis}
Other conditions cause muscle weakness in childhood and must be distinguished from DMD.
\begin{itemize}
    \item Becker muscular dystrophy, a milder form of the same genetic defect
    \item Other muscular dystrophies and congenital myopathies
    \item Spinal muscular atrophy, a nerve condition causing weakness
    \item Metabolic and mitochondrial myopathies
    \item Normal developmental delay without an underlying muscle condition
\end{itemize}

\section*{When to Seek Medical Care}
Any child with delayed motor milestones, unusual gait, or difficulty rising from the floor should be assessed by a GP and referred to a paediatrician. Early diagnosis allows families to access treatment and support that slow the disease and improve quality of life.

\textbf{Call 000 (or local emergency services) for:}
\begin{itemize}
    \item Sudden breathing difficulty or severe breathlessness
    \item Chest pain, palpitations, or fainting, which may indicate a heart problem
    \item Difficulty swallowing or choking
    \item Any sudden deterioration in a boy's condition
\end{itemize}

\section*{Diagnosis and Investigations}
Diagnosis of DMD is confirmed through a combination of tests.
\begin{itemize}
    \item \textbf{Blood tests:} creatine kinase levels are dramatically elevated because damaged muscle releases the enzyme
    \item \textbf{Genetic testing:} identifies the specific DMD gene mutation and confirms the diagnosis
    \item \textbf{Muscle biopsy:} rarely needed now that genetic testing is available, but may confirm absent dystrophin
    \item \textbf{Cardiac assessment:} regular echocardiograms and ECGs to monitor heart function
    \item \textbf{Respiratory assessment:} breathing tests to detect early respiratory weakness
    \item \textbf{Physiotherapy and occupational therapy assessment:} to plan support and equipment needs
\end{itemize}

\section*{Management}
Management of DMD is multidisciplinary and aims to slow progression, treat complications, and maintain quality of life.
\begin{itemize}
    \item \textbf{Corticosteroids:} medicines such as prednisolone and deflazacort slow the loss of muscle strength and extend the ability to walk
    \item \textbf{Physiotherapy:} stretching, positioning, and exercise programmes to maintain flexibility and prevent contractures
    \item \textbf{Occupational therapy:} splints, aids, and adaptations to help with daily activities
    \item \textbf{Respiratory support:} breathing exercises, cough assistance, and non-invasive ventilation when breathing weakness develops
    \item \textbf{Cardiac care:} regular monitoring and medication to support heart function
    \item \textbf{Surgery:} procedures to correct scoliosis or release contractures in some cases
    \item \textbf{Genetic therapies:} emerging treatments, including exon-skipping drugs and gene therapy, are becoming available for some mutations
    \item \textbf{Educational and psychological support:} tailored learning support and counselling for the child and family
\end{itemize}

\section*{Complications}
\begin{itemize}
    \item Loss of independent mobility in the early teenage years
    \item Scoliosis and joint contractures
    \item Respiratory failure from weakness of the breathing muscles
    \item Dilated cardiomyopathy and heart failure
    \item Swallowing difficulties and poor nutrition
    \item Learning and behavioural difficulties
    \item Emotional and practical strain on families
\end{itemize}

\section*{Prognosis and Outcomes}
Duchenne muscular dystrophy is a progressive condition without a cure, but the outlook has improved substantially over recent decades. With multidisciplinary care, most boys now maintain the ability to walk into their teenage years and live into their late twenties or thirties. Emerging genetic therapies offer the hope of slowing the disease further. The condition ultimately affects the heart and breathing muscles, and these systems require ongoing specialist monitoring and support.

\section*{Prevention and Self-Care}
\begin{itemize}
    \item Genetic counselling for families with a history of DMD, and carrier testing for at-risk women
    \item Prenatal testing and discussion of reproductive options where a mutation is known
    \item Early referral to a multidisciplinary neuromuscular service after diagnosis
    \item Maintaining an active daily stretching programme to delay contractures
    \item Staying up to date with cardiac and respiratory reviews
    \item Connecting with support organisations for information, equipment, and community
\end{itemize}

\section*{Sources and Further Reading}
The following reputable sources provide further information for healthcare students and patients seeking reliable, current advice about Duchenne muscular dystrophy.
\begin{itemize}
    \item Muscular Dystrophy Australia. \textit{Duchenne muscular dystrophy information and support.} https://www.mda.org.au/
    \item NHS. \textit{Duchenne muscular dystrophy: overview and management.} https://www.nhs.uk/conditions/duchenne-muscular-dystrophy/
    \item Mayo Clinic. \textit{Muscular dystrophy: symptoms and causes.} https://www.mayoclinic.org/diseases-conditions/muscular-dystrophy/
    \item MedlinePlus. \textit{Duchenne muscular dystrophy.} https://medlineplus.gov/genetics/condition/duchenne-muscular-dystrophy/
    \item Parent Project Muscular Dystrophy. \textit{Resources for families.} https://www.parentprojectmd.org/
    \item MSD Manual. \textit{Muscular dystrophies: professional overview.} https://www.msdmanuals.com/
\end{itemize}
